\begin{tbox}{\faLightbulb\quad تعريفات}
\begin{itemize}
    \item \textbf{المصفوفة} (جدول ثنائي البعد): جدول من الأسطر والأعمدة، نصل إلى كل خانة بفهرسين \code{matrix[row][col]}.
    \item \textbf{التصريح}: \code{int sales[3][4];} ثلاثة أسطر وأربعة أعمدة (\en{12} خانة).
    \item \textbf{التخزين}: الأسطر متتالية في الذاكرة (سطر بعد سطر).
    \item \textbf{المرور}: حلقتان متداخلتان؛ الخارجية على الأسطر والداخلية على الأعمدة (أو العكس لمجاميع الأعمدة).
\end{itemize}
\end{tbox}

\begin{center}
\begin{tikzpicture}[
    khana/.style={draw=brandTeal, fill=white, thick, minimum width=1.3cm,
        minimum height=0.75cm, font=\small\bfseries}]
    \foreach \row/\vals in {0/{120,95,140,180}, 1/{80,60,75,90}, 2/{200,210,190,250}} {
        \foreach \v [count=\col from 0] in \vals {
            \node[khana] at (-1.3*\col, -0.75*\row) {\en{\v}};
        }
        \node[font=\small, lawn=brandMuted] at (1.3, -0.75*\row) {\code{store \row}};
    }
    \foreach \col in {0,...,3}
        \node[font=\small, lawn=brandMuted] at (-1.3*\col, 0.65) {\code{day \col}};
    \draw[sahm, brandOrange] (-5.0,-1.5) -- (-4.5,-1.5);
    \node[anchor=east, font=\small] at (-5.0,-1.5) {\code{sales[2][...]}};
\end{tikzpicture}
\end{center}

\begin{center}
\begin{tabularx}{0.95\linewidth}{|c|X|}
\hline
\ras{العملية} & \ras{الفكرة} \\ \hline
مجموع سطر & الحلقة الداخلية على الأعمدة، والمجمِّع يُهيأ لكل سطر \\ \hline
مجموع عمود & نبدل الحلقتين: الخارجية على الأعمدة \\ \hline
المنقول & \code{transposed[col][row] = matrix[row][col]} \\ \hline
الجداء & \code{product[row][col] += left[row][mid] * right[mid][col]} (ثلاث حلقات) \\ \hline
البرمجة الديناميكية & كل خانة تُحسب من خانات سبق حسابها (فوق، يسار) \\ \hline
\end{tabularx}
\end{center}

\begin{obox}{\faExclamationTriangle\quad ترتيب الفهارس}
\code{matrix[row][col]}: السطر أولًا ثم العمود. الخلط بينهما يقرأ خانات خاطئة أو يتجاوز حدود المصفوفة.
\end{obox}

\insertstats{05}
